% SIFT 尺度空间构造示意：为减少层数取每组 s=2 个间隔（原文实验取 s=3），
% 每组 s+3=5 幅高斯图像、s+2=4 幅高斯差分图像；第二组边长减半。宽约 112mm。
\begin{tikzpicture}[x=1mm,y=1mm,font=\figurefont\small,
  gimg/.style={draw=FigureBlue,fill=FigureBlueFill,line width=.6pt},
  dimg/.style={draw=FigureYellow,fill=FigureYellowFill,line width=.6pt},
  link/.style={-{Stealth[length=1.8mm]},draw=FigureStroke,line width=.6pt},
  lab/.style={font=\figurefont\footnotesize,text=FigureMuted}]
  % octave 1: 5 gaussians stacked upward, size 16x10 slanted
  \foreach \k in {0,...,4}{
    \draw[gimg] (0+\k*1.5,\k*6) -- ++(16,0) -- ++(5,5) -- ++(-16,0) -- cycle;
  }
  \foreach \k in {0,...,3}{
    \draw[dimg] (40+\k*1.5,3+\k*6) -- ++(16,0) -- ++(5,5) -- ++(-16,0) -- cycle;
    \draw[link] (23+\k*1.5,5+\k*6) -- (39+\k*1.5,5.5+\k*6);
  }
  \node[lab,anchor=east] at (-1,2.5) {$\sigma$};
  \node[lab,anchor=east] at (-1,26.5) {$k^4\sigma$};
  \node at (13,36) {高斯图像 $L$};
  \node at (53,33) {高斯差分 $D$};
  \node[lab] at (31,-4) {相邻两层相减};
  % octave 2 (half size)
  \begin{scope}[yshift=-24mm]
    \foreach \k in {0,...,4}{
      \draw[gimg] (0+\k*.9,\k*3.4) -- ++(8,0) -- ++(2.5,2.5) -- ++(-8,0) -- cycle;
    }
    \foreach \k in {0,...,3}{
      \draw[dimg] (40+\k*.9,1.7+\k*3.4) -- ++(8,0) -- ++(2.5,2.5) -- ++(-8,0) -- cycle;
    }
    \node[lab,anchor=north west] at (-2,-1) {第二组：由尺度为 $k^s\sigma=2\sigma$ 的一层隔点下采样得到};
  \end{scope}
  \draw[link] (3,14.5) .. controls (-7,12) and (-7,-14) .. (0.2,-16);
  % 26-neighbour
  \begin{scope}[xshift=82mm,yshift=4mm]
    \foreach \l/\yy in {0/0,1/9,2/18}{
      \foreach \i in {0,1,2}{\foreach \j in {0,1,2}{
        \draw[FigureStroke,line width=.4pt,fill=white] (\i*3.6+\l*1.2,\yy+\j*2.2) rectangle ++(3.6,2.2);
      }}
    }
    \fill[FigureRed] (3.6+1.2+1.8,9+2.2+1.1) circle (0.9mm);
    \node[lab,anchor=west] at (13,1.5) {$D(\cdot,\cdot,\sigma_{i-1})$};
    \node[lab,anchor=west] at (13,10.5) {$D(\cdot,\cdot,\sigma_{i})$};
    \node[lab,anchor=west] at (13,19.5) {$D(\cdot,\cdot,\sigma_{i+1})$};
    \node at (8,29) {26 个邻点};
  \end{scope}
\end{tikzpicture}
